% !TeX root = ../../Book2.tex
\OptionalSection{极值问题}{b2:sec:2305}

算子有界以后，可以进一步研究最佳常数及达到它的输入。极值问题比有界性多出一个紧性要求；对称群常使普通的紧性失效。本节给出可核验的框架，不宣称所有曲面和指数都存在极值函数。

\subsection{最佳常数}
\label{b2:subsec:230501}

固定一个已证的界 \(E:L^2(S)\to L^q\)、\(q<\infty\)，令
\[
 \mathcal C=\sup_{\|g\|_2=1}\|Eg\|_q.
\]
当测度非零时 \(\mathcal C>0\)，有界性使其有限。若存在单位向量达到上确界，称为极值函数。最佳常数的数值、是否达到及所有极值函数的分类是三个不同问题，不应混成一个陈述。

\ProseParagraphBreak
即使算子范数有限，也未必有极值向量。例如 \(\ell^2\) 上的对角算子
\((x_n)\mapsto((1-1/n)x_n)\)（从 \(n=2\) 开始）范数为 \(1\)，但每个非零向量的范数都严格变小。单位坐标向量构成极值化序列，却没有强收敛子列。这个模型提醒我们必须寻找实际的紧性机制。

\subsection{极值化序列}
\label{b2:subsec:230502}

按上确界定义，可选 \(\|g_n\|_2=1\)、\(\|Eg_n\|_q\to\mathcal C\)。
Hilbert 空间弱紧性提供弱收敛子列，但弱极限可能为零。对于固定 \(x\)，曲面指数函数属于 \(L^2(S)\)，所以 \(g_n\rightharpoonup g\) 蕴含 \(Eg_n(x)\to Eg(x)\)；然而逐点收敛仍不保证 \(L^q\) 强收敛。

\ProseParagraphBreak
一个方便的充分判据是：若经对称归一化后 \(g_n\to g\) 在 \(L^2(S)\) 中强收敛，则连续性给出
\(\|g\|_2=1\)、\(\|Eg\|_q=\mathcal C\)。
该判据的证明很短，困难全在证明归一化后的强紧性，而不是上确界的形式定义。

\subsection{对称性}
\label{b2:subsec:230503}

球面的正交旋转保持输入及输出范数；乘一个常数相位也保持范数。环境平移由曲面调制实现，所以
\(g_n(\xi)=e^{-2\pi\mathrm ia_n\cdot\xi}g(\xi)\) 在 \(|a_n|\to\infty\) 时可使输出逃向无穷远。
在曲面局部放大极限中，还可能出现小帽集中，极限几何是切平面附近的抛物面；这不是单位球面精确的缩放对称，却是紧性分析中的真实损失机制。

\ProseParagraphBreak
归一化应分别处理精确对称与逼近尺度。除以相位、旋转、调制并不自动排除帽集中；排除帽集中也不自动排除多个远离的输出峰。每一种紧性损失都需要独立估计。

\subsection{极值函数}
\label{b2:subsec:230504}

若极值函数存在，在适当可微的范围 \(q>2\) 下，对变分 \(g+\epsilon h\) 求导可得到 Euler–Lagrange 方程
\[
 E^*(|Eg|^{q-2}Eg)=\mathcal C^qg
 \qquad(\|g\|_2=1).
\]
左侧按 \(L^{q'}\to L^2(S)\) 的伴随估计解释；因为
\(|Eg|^{q-2}Eg\in L^{q'}\)，该表达确实属于 \(L^2(S)\)。
求导时对复函数分别取实部方向和虚部方向，才能得到完整的复 Hilbert 空间等式。

\ProseParagraphBreak
这个方程是必要条件，不能反过来把任意解都称为全局极值函数。它还可能有局部临界点或鞍点；确认最佳常数仍需比较整个单位球。对于特殊曲面和特殊指数，卷积几何可能使常数与极值分类可显式计算，但本章不把这类特殊结论推广到一般情形。

\subsection{集中紧性}
\label{b2:subsec:230505}

集中紧性方法先寻找有正量的局部输出，再通过对称群把它移回固定区域；随后分离一个非零轮廓与余项，并用范数解耦判断余项能否继续携带极值量。\appref{b2:app:F}将说明消失、分裂和紧性三种模式，以及它们在变分问题中的作用。

\ProseParagraphBreak
一个典型的排除分裂原理是 \(q>2\) 时
\(a^{q/2}+(1-a)^{q/2}<1\)、\(0<a<1\)。
若两个轮廓在输出端充分分离、输入平方范数分别为 \(a\) 与 \(1-a\)，则两输出的 \(q\) 次范数之和至多
\(\mathcal C^q[a^{q/2}+(1-a)^{q/2}]\)，严格小于最优值。
使用这项原理之前必须证明轮廓解耦与没有其他损失；仅写这条数值不等式不能替代集中紧性的完整证明。

% 扩写定位：b2:subsec:230505
有一个可以完整证明的条件紧性判据，它把“输出没有逃走”与“输入达到极值”直接连接起来。设 \(S\) 紧致，\(E:L^2(S)\to L^q(\mathbb R^d)\) 有界，\(2<q<\infty\)，并有极值化序列 \(\|g_n\|_2=1\)。若在适当对称归一化后，
\[
 \lim_{R\to\infty}\sup_n\int_{|x|>R}|Eg_n(x)|^q\,\mathrm dx=0,
\]
则它具有强 \(L^2(S)\) 收敛子列，其极限为极值函数。

\ProseParagraphBreak

证明如下。取弱收敛子列 \(g_n\rightharpoonup g\)。由于 \(S\) 紧致，各阶导数都满足
\[
 |\partial_x^\alpha Eg_n(x)|
 \leq(2\pi)^{|\alpha|}
 \left(\int_S|\xi^\alpha|^2\,\mathrm d\sigma(\xi)\right)^{1/2}.
\]
因而输出在每个紧球上统一有界且等度连续。弱收敛给出逐点极限 \(Eg\)，Arzelà--Ascoli 及有限覆盖遂使 \(Eg_n\to Eg\) 在每个紧球上一致收敛。假设的统一尾部小性与 Fatou 控制极限尾部，故全空间 \(L^q\) 中也强收敛。于是 \(\|Eg\|_q=\mathcal C\)。弱下半连续性给出 \(\|g\|_2\leq1\)，而算子最佳界给出 \(\mathcal C\leq\mathcal C\|g\|_2\)；在非零算子情形 \(\mathcal C>0\)，所以 \(\|g\|_2=1\)。Hilbert 空间中弱收敛加范数收敛意味着强收敛，命题得证。

\ProseParagraphBreak

这一判据的假设很强，也明确指出了还需证明什么。调制产生的平移会破坏统一尾部；帽集中可能产生尺度越来越大的管，也可能破坏尾部。仅证明每个固定 \(n\) 的 \(Eg_n\) 可积，并不能把其尾部界变成对全部 \(n\) 一致的小量。因此上面的命题没有无条件断言一般曲面的极值函数存在，而是为具体问题提供一个可核对的紧性目标。
